\documentclass[11pt]{article}
\usepackage[margin=25mm]{geometry}
\usepackage{amsmath,amssymb,amsthm}
\usepackage[T1]{fontenc}
\usepackage[colorlinks=true,urlcolor=blue]{hyperref}
\newtheorem{theorem}{Theorem}
\setlength{\parindent}{0pt}\setlength{\parskip}{5pt}
\begin{document}
\begin{center}{\Large A Dirichlet polynomial with exact maximum N}\\[5pt]{\large Conjecture 00000001916}\\[5pt]zbw123456 \quad 14 October 2026\end{center}

We interpret the conjectured explicit factor as a constant independent of $N$, as required for the stated growth scale. An unrestricted $N$-dependent factor would impose no bound.
\begin{theorem}
For the admissible coefficients $a_n=1$,
\[
 D_N(t)=\sum_{n=1}^N n^{it},\qquad \sup_{t\in\mathbb R}|D_N(t)|=N.
\]
This maximum is not $O\bigl(N^{1/2}(\log\log N)^{1/4}\bigr)$ as $N\to\infty$.
\end{theorem}
\begin{proof}
Every coefficient has modulus one. For real $t$, $n^{it}=\exp(it\log n)$ has modulus one, so the triangle inequality gives $|D_N(t)|\le N$. At $t=0$ each term is one, and equality holds. Thus the supremum is attained and is exactly $N$.

For sufficiently large real $x$, $\log x\ge1$ and
\[
0\le(\log\log x)^{1/4}\le(\log x)^{1/4}=o(x^{1/2}).
\]
Multiplication by $x^{1/2}$ shows that the claimed scale is $o(x)$. Restricting to positive integers preserves this relation. If $N$ were bounded by a constant times the claimed scale, transitivity would give $N=o(N)$, which is impossible: its defining estimate with constant $1/2$ would give $N\le N/2$ for all sufficiently large positive integers.
\end{proof}
\textbf{Lean correspondence.} \texttt{dirichlet} is the actual finite complex exponential sum, indexed by $n+1$ to avoid the logarithm of zero. \texttt{maximum\_exact} proves the \texttt{IsGreatest} property for the range of its norm, so no informal supremum computation is assumed. \texttt{scale\_littleO} uses the proved logarithm--power asymptotics of Mathlib, and \texttt{no\_claimed\_bound} negates the exact big-$O$ assertion. The witness is deterministic, as in the question; no random-coefficient hypothesis appears in its text.


\textbf{Verification.} Lean 4.19.0 with Mathlib pinned to
\texttt{c44e0c8ee63ca166450922a373c7409c5d26b00b}.
The accompanying project and verification transcript give the exact theorem and axiom audit. No proof placeholders, native evaluator proof bypass, or custom axiom declarations are used. This is a prepared submission; official acceptance has not been claimed.
\end{document}
